% ============================================================================
% CHAPTER 5 TEACHER'S MANUAL
% Complete instructional guide for teaching HITL intervention design
% ============================================================================
% Source: /markdown/ch5/Ch5T_updated.md
% Note: This file is for the Instructor's Edition, not the main book
% ============================================================================

\chapter{Teacher's Manual: Chapter 5}
\label{ch:teacher-ch05}

\begin{chapterquote}{Complete instructional guide}
Teaching the art of designing human-AI interactions using \emph{The Art of Asking for Help Without Being Annoying}.
\end{chapterquote}

% ============================================================================
\section{Course Integration Guide}
% ============================================================================

\subsection{Learning Objectives}

By the end of this chapter, students will be able to:

\paragraph{Knowledge (Remembering \& Understanding).}
\begin{itemize}
    \item Name the three primary design levers for HITL interventions: Timing, Frequency, Format (T-F-F)
    \item Define the trust-attention-response triangle and describe how each vertex can fail
    \item Explain the ``Lazy User Assumption'' and why it leads to bad intervention design
    \item Describe how question framing affects response quality, using the 911 dispatcher example
    \item Identify the core difference between a well-timed and a well-formatted ask
\end{itemize}

\paragraph{Application \& Analysis.}
\begin{itemize}
    \item Evaluate existing alert designs using the T-F-F framework and trust-attention-response triangle
    \item Identify when poor alert performance is due to timing, frequency, or format failures
    \item Apply cognitive load principles to assess query design choices
    \item Analyze fraud alert, 2FA, and content moderation examples using the chapter's framework
\end{itemize}

\paragraph{Synthesis \& Evaluation.}
\begin{itemize}
    \item Design HITL interventions that respect timing, manage frequency, and optimize format
    \item Critique real alert systems and propose specific, implementable improvements
    \item Argue for specific design choices using trust-attention-response reasoning
    \item Evaluate the information value of different query formats for a specific domain
\end{itemize}

\subsection{Prerequisites}
\begin{itemize}
    \item \textbf{Chapters 1--4 (required):} Five Dimensions framework, uncertainty types, calibration, threshold decisions; Chapter~5 is where the routing decision (Ch.\ 4) meets the human moment
    \item \textbf{No technical background required} for the main chapter
    \item \textbf{Appendix~5A} uses information theory (entropy, mutual information) and basic probability
\end{itemize}

\subsection{Course Positioning}

\begin{description}
    \item[Standard sequence] After Chapter~4 (threshold decisions) and before the domain-specific chapters --- Chapter~5 addresses what happens at the moment of human-AI handoff
    \item[HCI emphasis] Core content on interactive alert design; connections to notification research, dual-process theory
    \item[ML pipeline courses] The feedback integration argument shows how question format determines training signal quality
\end{description}

% ============================================================================
\section{Key Concepts for Instruction}
% ============================================================================

\subsection{The Central Design Insight}

The chapter's core argument: \textbf{the moment of the ask is a design problem, not just an engineering trigger}. Even when detection, routing, and threshold decisions are correct, the ask can fail if its design does not respect human psychology. One word --- ``personally'' --- created a 40\% improvement in the Louisville 911 dispatcher study. This is not a marginal efficiency gain; it is the difference between a system that works and one that does not.

\begin{technote}[Teaching Tip]
Open with the 911 dispatcher story before introducing any vocabulary. Ask students: ``What would you change about the question `Is there someone there who needs help?''' Then reveal the answer (one word: ``personally''; and ``right now''). The concreteness of one word creating 40\% improvement is vivid. Then ask: ``Why did one word matter so much?'' Lead students to: the question's framing changes what the caller searches for in memory, and therefore what information they report.
\end{technote}

\subsection{The Three Design Levers: T-F-F}

\begin{table}[htbp]
\caption{The three HITL design levers: definitions and failure examples}
\label{tab:tff-ch05}
\centering
\begin{tabular}{@{}p{2cm}p{4.5cm}p{4.5cm}@{}}
\toprule
\textbf{Lever} & \textbf{Definition} & \textbf{Failure Example} \\
\midrule
Timing & When the ask is sent & Fraud alert at 3 AM for a morning charge; context has evaporated \\
Frequency & How often asks are sent & iPhone update prompt: dismissed after repetition, even for critical updates \\
Format & How the ask is phrased and presented & ``Your account has been flagged'' vs.\ ``Did you make this \$47.23 purchase at Costa Coffee at 2:15 PM today?'' \\
\bottomrule
\end{tabular}
\end{table}

\begin{keyconcept}[Mnemonic: T-F-F]
Help students remember: \textbf{T}iming = when you ask; \textbf{F}requency = how often you ask; \textbf{F}ormat = how you ask. These are not independent --- a perfect format at the wrong time fails. All three must be optimized together.
\end{keyconcept}

\subsection{The Trust-Attention-Response Triangle}

Draw this on the board at the start of the relevant section:

\begin{table}[htbp]
\caption{Trust-attention-response triangle: failure modes at each vertex}
\label{tab:triangle-ch05}
\centering
\begin{tabular}{@{}p{2.5cm}p{4.5cm}p{4cm}@{}}
\toprule
\textbf{Vertex} & \textbf{Failure Mode} & \textbf{Example} \\
\midrule
Trust & Alerts dismissed because precision is too low & Clinical alerts: physicians override 90\%+ when precision is low \\
Attention & Alert arrives when reviewer cannot engage & Interrupt mid-complex task; 3 AM fraud alert \\
Response & Question format gets useless answers & ``Was this transaction legitimate?'' when user can't remember \\
\bottomrule
\end{tabular}
\end{table}

\subsection{The Lazy User Assumption}

State the assumption explicitly: most bad alert design implicitly assumes ``if users don't engage, it is because they are lazy.'' Then attack it:
\begin{enumerate}
    \item Empirically wrong --- research shows users engage when design respects their time
    \item Fails on its own terms --- friction does not select for motivated users; it just loses them
    \item The right assumption: ``users will engage if the ask is worth their time''
\end{enumerate}

\textbf{Design implication:} What makes an ask worth a user's time? It should be: (1) rare (surprising); (2) timely (user can respond with current context); (3) clear (user understands what is being asked); (4) consequential (user can see that their response changes something).

% ============================================================================
\section{Lecture Planning Guide}
% ============================================================================

\subsection{50-Minute Lecture Structure}

\subsubsection{Opening: The 911 Story (7 minutes)}

Tell the Louisville story slowly. Before revealing the improvement, ask students: ``What would you change about the question `Is there someone there who needs help?''' Then reveal the answer.

\begin{trythis}
\textbf{Key question:} ``Why did one word matter so much?'' Lead students to: the question's framing changes what the caller searches for in memory. ``Personally'' and ``right now'' anchor the caller to their immediate sensory experience rather than a general assessment.
\end{trythis}

\subsubsection{Part 1: The Three Levers (10 minutes)}

Walk through T-F-F with vivid examples:
\begin{itemize}
    \item \textbf{Timing failure:} Fraud alert at 3 AM. Who gets these at 3 AM? Why does it happen? (Batch processing, not real-time.) What does a half-awake approval actually contribute?
    \item \textbf{Frequency failure:} The iPhone update prompt. How many have dismissed without reading? Why? What does that mean for the rare important update?
    \item \textbf{Format failure:} ``Your account has been flagged for unusual activity'' vs.\ ``Did you make this \$47.23 purchase at Costa Coffee at 2:15 PM today? Reply YES or NO.''
\end{itemize}

\paragraph{Key discussion question.} Why did the format failure example feel intuitively worse than the timing failure? (Because format failure destroys the information value of the response --- it gets you a useless answer. Timing failure at least gets the right question at the wrong time.)

\subsubsection{Part 2: The Trust-Attention-Response Triangle (12 minutes)}

Introduce each vertex with a failure example (Table~\ref{tab:triangle-ch05}).

\begin{trythis}
\textbf{Class discussion (3 min):} ``Which vertex of the triangle is hardest to recover from once lost? Which is most under the designer's control?''\\[0.5em]
\textbf{Strong answer:} Trust is hardest to recover (erodes fast, builds slowly). Format is most under the designer's control (each question is a deliberate choice).
\end{trythis}

\subsubsection{Part 3: The Lazy User Assumption (8 minutes)}

State the assumption, attack it, then derive the design implication. What makes an ask worth a user's time?

\subsubsection{Part 4: Case Studies (8 minutes)}

Comparative analysis of 2FA and fraud alerts:

\begin{description}
    \item[2FA: Bad] Push notification, no context, tap to approve --- the MFA fatigue attack exploits reflexive approval
    \item[2FA: Good] Number matching, context about device/location, prominent ``I didn't initiate this'' option
    \item[Fraud: Bad] Generic flagging, delayed, requires calling bank
    \item[Fraud: Good] Specific merchant/amount/time, immediate, binary reply
\end{description}

Tie both back to T-F-F: identify which lever is well-designed in each case.

\subsubsection{Wrap-Up: Feedback Integration (5 minutes)}

The ask is useless if the response disappears. The question format determines what information is captured. The feedback loop determines whether that information improves the system.

Close the five-chapter arc: detection (Ch.\ 2) $\to$ learning (Ch.\ 3) $\to$ thresholds (Ch.\ 4) $\to$ asking (Ch.\ 5). Each chapter has answered a successively more specific question about the HITL moment.

\begin{table}[htbp]
\caption{50-minute lesson plan for Chapter~5}
\label{tab:lesson-plan-ch05}
\centering
\begin{tabular}{@{}p{2.5cm}p{6cm}p{4cm}@{}}
\toprule
\textbf{Time} & \textbf{Activity} & \textbf{Materials} \\
\midrule
0:00--0:07 & Opening: 911 story + discussion & Ch.\ 5 opening narrative \\
0:07--0:17 & Three levers (T-F-F) with examples & Slides with real alert comparisons \\
0:17--0:29 & Trust-attention-response triangle + discussion & Whiteboard triangle diagram \\
0:29--0:37 & Lazy User Assumption + correction & Discussion \\
0:37--0:44 & Case studies: 2FA and fraud alerts & Comparative slides \\
0:44--0:49 & Feedback integration + chapter arc close & \\
0:49--0:50 & Assign Try This; preview domain chapters & \\
\bottomrule
\end{tabular}
\end{table}

\subsection{75-Minute Extended Session}

\paragraph{Add: Alert Triage Activity (15 min).}
See Activity~1 below. Students rank six alert designs and explain their ranking.

\paragraph{Add: HITL Redesign Workshop (10 min).}
See Activity~2 below. Small groups redesign a specified intervention system.

% ============================================================================
\section{Discussion Questions by Level}
% ============================================================================

\subsection{Introductory Level}

\begin{enumerate}
    \item \textbf{Recognition:} ``The 911 dispatcher story shows that the wording of a question can dramatically change the response. Give two examples from your own experience where how you were asked a question changed how you answered it.''\\
    \textit{Target: Strong answers identify specific episodes where phrasing changed what the respondent searched for. Example: ``Do you have any questions?'' vs.\ ``What's one thing that's still unclear?'' --- the latter produces genuine reflection.}

    \item \textbf{Analysis:} ``A mobile banking app sends you: `Unusual activity detected. Open the app to review.' Apply the T-F-F framework: what does this notification do well and do poorly?''\\
    \textit{Target: Timing (likely unpredictable, context may be stale), Frequency (if sent for many events, users learn to dismiss), Format (no actionable specifics; requires 3--4 steps before user can even assess). Poor on all three.}

    \item \textbf{Application:} ``Redesign the banking notification from Question~2 to be better on all three dimensions. What information would you include? When would you send it? How would you prevent alert fatigue?''\\
    \textit{Target: ``Did you make a \$127.43 purchase at GameStop on Friday at 6:15 PM? Reply YES to confirm or NO to freeze your card.'' Sent within 60 seconds; only for statistically anomalous transactions; binary reply with stated consequence.}

    \item \textbf{Comparison:} ``Compare two-factor authentication with push notifications (tap to approve) vs.\ number matching. Which is better and why? Use the trust-attention-response triangle.''\\
    \textit{Target: Push fails on response dimension --- approval requires no information from user that an attacker couldn't provide. Number matching requires active engagement with information an attacker cannot have. Both score equally on timing and frequency; difference is entirely format.}
\end{enumerate}

\subsection{Intermediate Level}

\begin{enumerate}
    \item \textbf{Alert Fatigue:} ``A hospital's clinical decision support system generates 500 alerts per doctor per day, and doctors override 95\% of them. Diagnose this system using the trust-attention-response triangle. Which vertex is failing most critically? Propose three specific interventions.''\\
    \textit{Target: Trust failure (5\% precision). Also attention failure (one alert every 57 seconds). Three interventions: precision triage (eliminate low-actionability categories), adaptive thresholds, batching + contextual presentation at workflow breaks.}

    \item \textbf{Design Challenge:} ``Design a HITL intervention for a plagiarism detection system used by university professors. Specify: what gets flagged, when the professor is notified, what format the notification takes, and how the professor's response feeds back to the system.''\\
    \textit{Target: Trigger on high similarity scores (e.g., >0.75) in uncertain zone; notify after submission batch closes (not mid-grading flow --- timing matters); format includes similarity percentage + specific source + excerpt comparison; professor response (confirm/dismiss/escalate) feeds into model recalibration.}

    \item \textbf{Information Value:} ``A fraud alert system can send a specific alert (`Did you make this \$23.45 purchase at Starbucks at 9:14 AM?') or a generic alert (`A transaction may need your review'). Why is the specific alert better even though both ask the same underlying question?''\\
    \textit{Target: Specific alert activates episodic memory. User can directly check their experience of that morning. Generic alert activates only general assessment: no new information. Information gain (entropy reduction) from specific alert is high; from generic alert is near zero.}
\end{enumerate}

\subsection{Advanced Level}

\begin{enumerate}
    \item \textbf{System Design:} ``A content moderation platform employs 200 moderators reviewing 50,000 pieces of flagged content per day. Using all five dimensions of the HITL framework, design a complete intervention system including: routing criteria, interface design, cognitive load management, alert frequency management, and feedback integration.''

    \item \textbf{Research Design:} ``Design an experiment to measure the information value of different fraud alert formats. What would you measure? How would you quantify `information quality' from a human response? What confounds would you control for?''

    \item \textbf{Dynamic Optimization:} ``Using the RL formulation from Appendix~5A, explain why the optimal HITL policy should be context-dependent (different thresholds for different users, times, and situations) rather than a fixed rule. What data would you need to learn such a policy?''
\end{enumerate}

% ============================================================================
\section{Hands-On Activities}
% ============================================================================

\subsection{Activity 1: Alert Triage (15 minutes)}

\textbf{Setup:} Present six alert designs (Table~\ref{tab:alert-triage}). Students rank them from best to worst and explain their ranking.

\begin{table}[htbp]
\caption{Six alerts for triage ranking exercise}
\label{tab:alert-triage}
\centering
\begin{tabular}{@{}p{1cm}p{7cm}p{3.5cm}@{}}
\toprule
\textbf{\#} & \textbf{Alert Text / Description} & \textbf{Assessment} \\
\midrule
1 & ``Your account requires your attention'' & Poor all dimensions \\
2 & ``Did you make this \$237.00 purchase at Best Buy at 4:47 PM today? Reply YES or NO'' & Good all dimensions \\
3 & Push notification: ``New login detected --- Tap to approve'' & Good timing, poor format \\
4 & Email 6 hours later: ``We noticed a large purchase on your account'' & Poor timing, poor format \\
5 & Phone call: ``We're calling to verify a recent transaction\ldots'' & Good format, bad attention \\
6 & In-app banner: ``Review 14 flagged items'' & Timing/format problems \\
\bottomrule
\end{tabular}
\end{table}

\textbf{Debrief:} The reveal is that Alert~2 should rank highest and Alert~1 lowest. More interesting: Alert~3 (good timing, bad format) vs.\ Alert~4 (bad timing, bad format) --- both are bad, but why differently? Alert~5 (good format, bad attention) shows that format alone is insufficient.

\subsection{Activity 2: HITL Intervention Redesign (20 minutes)}

\textbf{Setup:} Groups of 3--4 redesign the intervention system for a specific scenario.

\paragraph{Scenarios.}
\begin{description}
    \item[A: Premium credit card fraud] International customer, high-value transaction in unfamiliar city, unusual time
    \item[B: University plagiarism detection] Professor reviewing 30 student submissions; one flagged at 78\% similarity
    \item[C: Smart home security] Door left unlocked when no one is home; AI is 82\% confident
    \item[D: Medical documentation AI] Physician dictation transcription detected possible error on drug dosage
\end{description}

\textbf{Deliverable:} Each group specifies: trigger condition, notification text, timing, frequency controls, response format, and feedback integration mechanism.

\subsection{Activity 3: Try This Debrief (10 minutes)}

\textbf{Second-session activity.} Students share their audit of alerts they experienced between classes.

\paragraph{Discussion prompts.}
\begin{itemize}
    \item Which was the worst-timed alert you received?
    \item Which format felt most respectful of your time?
    \item Did any alert feel like it was designed to extract your specific knowledge vs.\ just get an approval?
    \item Which alert taught you the most about good design by being bad?
\end{itemize}

% ============================================================================
\section{Common Student Misconceptions}
% ============================================================================

\subsection{Misconception 1: ``More Information in the Ask = Better Ask''}

\paragraph{Student thinking:} ``A fraud alert should give you everything: merchant, amount, time, location, device, and previous purchase history.''

\paragraph{Correction strategy.}
\begin{itemize}
    \item Cognitive load theory shows accuracy peaks when complexity matches working memory capacity
    \item A comprehensive but overwhelming ask produces lower-quality responses than a well-focused simpler one
    \item Good design finds the minimum sufficient information, not the maximum available information
    \item The goal is to activate the right memory at the right level of specificity
\end{itemize}

\subsection{Misconception 2: ``Automated Confirmation Is Fine for Low-Stakes Asks''}

\paragraph{Student thinking:} ``For routine things, it doesn't matter if people approve without reading.''

\paragraph{Correction strategy.}
\begin{itemize}
    \item The MFA fatigue attack exploits exactly this logic --- users who habitually approve without reading become attack vectors
    \item Habitual non-reading transfers across contexts: users who learn to dismiss ``low-stakes'' alerts apply the same behavior to higher-stakes ones (trust erosion across contexts)
    \item The categorization of ``low-stakes'' is itself a design decision that can be wrong
\end{itemize}

\subsection{Misconception 3: ``The Format Problem Is Just a UX Issue''}

\paragraph{Student thinking:} ``This is about making things look nice, not about fundamental AI design.''

\paragraph{Correction strategy.}
\begin{itemize}
    \item The information the human provides is the primary input to the system's learning loop
    \item A badly formatted question produces a response with less useful information for model improvement
    \item This is not UX polish --- it is core ML pipeline engineering: garbage format = garbage training signal
    \item The information-theoretic argument: question format determines entropy reduction in the response
\end{itemize}

\subsection{Misconception 4: ``Alert Fatigue Is an Individual User Problem''}

\paragraph{Student thinking:} ``Some users are just lazy or have too many notifications.''

\paragraph{Correction strategy.}
\begin{itemize}
    \item Alert fatigue is a property of system design, not user character
    \item Well-designed systems with sufficient precision maintain engagement
    \item The dynamic trust model shows that any system with low alert precision will degrade trust over time, regardless of user motivation
    \item The clinical alert fatigue literature documents this in among the most dedicated professionals: physicians
\end{itemize}

% ============================================================================
\section{Assessment Options}
% ============================================================================

\subsection{Option 1: Alert Audit Report (500--750 words)}

\textbf{Prompt:} ``Audit five alerts or HITL interventions you encountered this week. For each: (a) identify which of timing, frequency, format is best and worst; (b) apply the trust-attention-response triangle; (c) propose one specific improvement. Conclude with a pattern you noticed across the five.''

\paragraph{Rubric.}
\begin{itemize}
    \item \textbf{T-F-F application (35\%):} Correctly identifies and distinguishes the three levers for each alert
    \item \textbf{Triangle analysis (25\%):} Identifies which vertex is failing and why
    \item \textbf{Improvement quality (25\%):} Specific, feasible, addresses the identified failure mode
    \item \textbf{Pattern identification (15\%):} Cross-alert pattern is genuine and well-supported
\end{itemize}

\subsection{Option 2: HITL Design Document}

\textbf{Prompt:} ``Design the complete human-AI collaboration interface for a domain of your choice. Include: trigger conditions, notification design (with mockup or detailed description), frequency controls, response options, and feedback integration mechanism. Justify each design decision using chapter concepts.''

\paragraph{Rubric.}
\begin{itemize}
    \item \textbf{Completeness (30\%):} All six required components present
    \item \textbf{Conceptual grounding (35\%):} Each design decision justified using T-F-F, trust-attention-response, or Lazy User Assumption
    \item \textbf{Feedback integration (20\%):} Explicit mechanism; connected to model improvement
    \item \textbf{Feasibility (15\%):} Design is actually implementable
\end{itemize}

\subsection{Option 3: Case Analysis}

\textbf{Prompt:} ``The chapter describes how professional content moderators' decision quality is affected by interface design, throughput pressure, and context availability. Analyze a content moderation platform you are familiar with (or based on published research). What would you change?''

% ============================================================================
\section{Connections to Other Chapters}
% ============================================================================

\subsection{Backward Links}
\begin{itemize}
    \item \textbf{Chapter~1:} Alert fatigue first introduced; the Goldilocks problem; Nest thermostat as asking badly
    \item \textbf{Chapter~2:} Uncertainty signal must be surfaced before it can be asked about; calibration determines whether the ask is triggered at the right moment
    \item \textbf{Chapter~3:} The training signal from human responses depends on the quality of the question --- bad format = noisy training data
    \item \textbf{Chapter~4:} The HITL band defines when asking happens; format determines what is learned from each ask
\end{itemize}

\subsection{Forward Links}
\begin{itemize}
    \item \textbf{Chapter~6:} Interface design expands on Format in depth --- the visual and cognitive design of the ask itself
    \item \textbf{Chapter~7:} Annotation science is structured asking at scale --- every principle in Chapter~5 applies to annotation guideline design
    \item \textbf{Domain chapters (Part~II):} Each domain applies the T-F-F framework in a specific professional context
\end{itemize}

\bigskip
\begin{center}
\rule{0.5\textwidth}{0.4pt}
\end{center}
\bigskip

\noindent\textit{This teacher's manual provides everything needed to teach Chapter~5 on HITL intervention design. The central insight --- that the ask is a design problem, not just an engineering trigger --- closes the first arc of the book and opens the interface design arc that follows.}

\medskip
\noindent\textit{Next: Chapter~6 deepens the interface question: not just how to ask, but how to present AI output to humans in ways that support rather than replace their judgment.}
